\section{A degree-5 generator in 44 variables}
\label{sec:deg5}

\subsection{Starting point}

We start from the system of \cite{jsww1976}, which we reproduce in the form our
tooling uses. It has $14$ equations in the parameter $k$ and $25$ unknowns
$a,b,\dots,z$; the primes are the values $k+2$ for which it is solvable, and by
Proposition~\ref{prop:gen} it yields a generator with
$\nu=26$ and $\delta=1+2\cdot 12=25$.

\subsection{The construction}

The generator is obtained in two steps.

\paragraph{Flattening to degree $2$.}
Twenty subexpressions are named, one new unknown each. The choice of names is
the output of an optimizer: a catalogue of candidate subexpressions is built
from the syntactic tree of the system, a selection is encoded as a
satisfiability problem, and an SMT solver is asked for a selection of minimum
size such that every equation reduces to degree at most~$2$. We stress at once,
and repeat in Section~\ref{sec:limits}, that the resulting optimality statement
is \emph{relative to that encoding}: it says that no smaller selection exists
\emph{within the catalogue}, not that no smaller flattening exists.

\paragraph{Two eliminations.}
On the flattened system, two unknowns are determined linearly with nonnegative
right-hand sides,
\[
  q = h+j+wz, \qquad y = l+n+v ,
\]
so Remark~\ref{rem:crit} applies and both may be eliminated with no auxiliary
argument. No third elimination is available: $q$ and $e$ are determined by the
\emph{same} equation, so they compete for it, and eliminating $z$ raises the
degree above~$2$. We measured all six orders.

\paragraph{The count.}
\[
  44 \;=\; \underbrace{26}_{\text{\cite{jsww1976}}}
        \;+\; \underbrace{20}_{\text{names}}
        \;-\; \underbrace{2}_{\text{eliminations}} .
\]

\begin{theorem}
\label{thm:deg5}
The system displayed in \S\ref{ssec:flat-system}, in $44$ variables
and with every equation of degree at most $2$, is equisatisfiable over $\N$ with
the system of \textup{\cite{jsww1976}}. Consequently
\[
  Q_{44} \;=\; (k+2)\Bigl(1-\textstyle\sum_i P_i^2\Bigr)
\]
is a generator of the primes with $\nu=44$ and $\delta=5$.
\end{theorem}

Both directions of the equisatisfiability are machine-checked; see
Section~\ref{sec:lean}.

\subsection{The system}
\label{ssec:flat-system}

Write $\mu_1,\dots,\mu_{20}$ for the named subexpressions. The twelve principal
equations are
%% GENERADO por paper/generar_sistemas.py — no editar a mano
%% 44 variables, 12 ecuaciones principales + 20 definitorias
\begin{align}
  j + 2h + h k + h \mu_{1} + j k + j \mu_{1} &= z,\\
  h + j + p + z + 2n + w z &= e,\\
  1 + 16\mu_{4} + 16\mu_{3} \mu_{4} + 32\mu_{4} n &= f^{2},\\
  1 + \mu_{7} + \mu_{6} \mu_{7} + 2a \mu_{7} &= o^{2},\\
  1 + \mu_{6} \mu_{8} &= l^{2} + n^{2} + v^{2} + x^{2} + 2l n + 2l v + 2n v,\\
  1 + 16\mu_{10}^{2} &= u^{2} + 16\mu_{9}^{2},\\
  1 + \mu_{14} \mu_{15} &= \mu_{15} + \mu_{12}^{2} + x^{2} + 2\mu_{12} x,\\
  1 + \mu_{16} \mu_{6} &= l^{2} + m^{2},\\
  1 + k + a i &= i + l,\\
  p + a l + 2a b + 2a \mu_{17} &= l + m + 2b + l n + \mu_{17} n + 2b n,\\
  h + j + a l + a n + a v + w z + 2a s + 2\mu_{19} s &= l + n + v + x + 2s + l p + \mu_{18} s + n p + p v + 2p s,\\
  z + l \mu_{19} + 2\mu_{19} t &= t + l \mu_{18} + m p + \mu_{18} t.
\end{align}

\begin{align}
  \mu_{1} &= 2g + g k,\\
  \mu_{2} &= k^{2},\\
  \mu_{3} &= n^{2},\\
  \mu_{4} &= 2 + \mu_{2}^{2} + 7k + 9k^{2} + 5k \mu_{2},\\
  \mu_{5} &= e^{2},\\
  \mu_{6} &= a^{2},\\
  \mu_{7} &= \mu_{5}^{2} + 2e \mu_{5},\\
  \mu_{8} &= l^{2} + n^{2} + v^{2} + 2l n + 2l v + 2n v,\\
  \mu_{9} &= \mu_{8} r,\\
  \mu_{10} &= a \mu_{9},\\
  \mu_{11} &= u^{2},\\
  \mu_{12} &= c u,\\
  \mu_{13} &= n + 4d l + 4d n + 4d v,\\
  \mu_{14} &= \mu_{20}^{2},\\
  \mu_{15} &= \mu_{13}^{2},\\
  \mu_{16} &= l^{2},\\
  \mu_{17} &= b n,\\
  \mu_{18} &= p^{2},\\
  \mu_{19} &= a p,\\
  \mu_{20} &= a + \mu_{11}^{2} - a \mu_{11}.
\end{align}



\noindent
Every equation has degree at most~$2$, so $\delta=1+2\cdot 2=5$. The variables
are the parameter $k$, the $23$ surviving unknowns
$a,b,c,d,e,f,g,h,i,j,l,m,n,o,p,r,s,t,u,v,w,x,z$, and the $20$ names, for $44$ in
total. The \LaTeX{} of the display above is generated mechanically from the
same file the Lean kernel checks, so it cannot drift from the verified object.

\subsection{The twentieth name}

Nineteen of the twenty witnesses exhibited by the completeness direction are
sums and products of nonnegative quantities and therefore lie in $\N$ for
syntactic reasons. One does not:
\[
  \mu_{20} \;=\; a+u^2(u^2-a) .
\]
This matters, and in a way that is easy to miss. If $\mu_{20}$ could be
negative, the flattened system would have \emph{no} solution for the prime in
question, so the generator would fail to emit it. That is a failure of
completeness, not of soundness, and it is silent: the pair $(44,5)$ would still
be produced.

\begin{proposition}
\label{prop:mu20}
Let $a,r,y,u\in\Z$ with $a\ge 2$, $r,y\ge 0$ satisfy equation~\textup{(7)} of
the system, namely $16r^2y^4(a^2-1)+1=u^2$. Then $a+u^2(u^2-a)\ge 1$.
\end{proposition}

\begin{proof}
Put $t=u^2$. The identity
\[
  a+t(t-a)-1 \;=\; (t-1)(t-a+1)
\]
reduces the claim to $(t-1)(t-a+1)\ge 0$. By hypothesis
$t-1=16s(a^2-1)$ with $s=r^2y^4\ge 0$. If $s=0$ then $t=1$ and the first factor
vanishes. If $s\ge 1$ then, since $a\ge 2$ gives $a^2-1\ge 3$,
\[
  t-1 \;=\; 16s(a^2-1)\;\ge\; 16(a^2-1)\;\ge\; a-1 ,
\]
the last step because $16a^2-a-15\ge 0$ for $a\ge 1$; hence $t\ge a$ and both
factors are nonnegative.
\end{proof}

The hypothesis $a\ge 2$ is not an assumption: it is proved from five of the
fourteen equations, independently of what the system represents
(Section~\ref{sec:lean}). Proposition~\ref{prop:mu20} is formalized in Lean;
before it was proved, this bound had only been checked by exhaustive search over
the $2569$ triples $(a,r,y)$ with $a<45$, $r,y<30$ satisfying~(7), which is not
a proof.

\subsection{What we could not improve}

The pair $(44,5)$ is the best we obtained, and we report the attempts that
failed because they delimit where an improvement would have to come from.

\begin{center}
\begin{tabular}{@{}ll@{}}
\toprule
\textbf{Attempt} & \textbf{Outcome} \\
\midrule
More eliminations & Structurally capped at $2$ (see above); all six orders measured \\
Naming degree-$1$ subexpressions & Yields $+1$ name and $+1$ elimination: no change \\
Enlarging the catalogue & $465\to946\to3469$ candidates; optimum stays at $20$ \\
Drop-one test & Removing any of the $20$ names breaks materialization \\
Aligning the partition rules & Encoding was stricter than the materializer; optimum unmoved \\
The quotient-method system & Unsatisfiable under the structural filter; $(56,5)$ without it \\
The stronger Pell reparametrization & Gives $49$ variables, worse \\
\bottomrule
\end{tabular}
\end{center}

None of this shows that $44$ is optimal. It shows that seven distinct attacks,
each of which we expected to help, did not.
